\subsection{与等价关系相容的关系}

让 \(\mathbf{R}\{x, x'\}\) 是一个等价关系，且设 \(\mathbf{P}\{x\}\) 是一个关系。关系 \(\mathbf{P}\{x\}\) 称为与等价关系相容 \(\mathbf{R}\{x, x'\}\) （关于 \(x\) ）如果我们有

\[
(\mathrm{P} \{x \} \text {   且   } \mathrm{R} \{x, y \}) \Longrightarrow \mathrm{P} \{y \},
\]

其中 \(y\) 表示一个既不出现在 \(P\) 中也不出现在 \(R\) 中的字母。

中的字母。例如，由C43（第一章，§5，第1号）可知，每个关系 \(\mathbf{P}\{x\}\) 与等价关系相容 \(x = x'\) .

C56. 设 \(R\{x, x'\}\) 是集合 E 上的一个等价关系，并设 \(P\{x\}\) 是一个不含字母 \(x'\) 且（关于 x）与等价关系 \(R\{x, x'\}\) 相容的关系。那么，如果 t 不出现在 \(P\{x\}\) 中，关系“ \(t \in E/R\) 且 \((\exists x)(x \in t \text{ 且 } P\{x\})\) ”等价于关系“ \(t \in E/R\) 且 \((\forall x)((x \in t) \Longrightarrow P\{x\})\) ”。

设 \(t \in \mathbb{E} / \mathbb{R}\) 。如果存在 \(a \in t\) 使得 \(\mathbb{P}\{a\}\) ，那么对每个 \(x \in t\) 都有 \(\mathbb{R}\{a, x\}\) ，因而有 \(\mathbb{P}\{x\}\) 。所以 \((\exists x)(x \in t \text{ 且 } \mathbb{P}\{x\})\) 蕴含 \((\forall x)((x \in t) \Longrightarrow \mathbb{P}\{x\})\) 。逆命题显然成立，因为 \(t \in \mathbb{E} / \mathbb{R}\) 蕴含 \(t \neq \emptyset\) .

关系

\[
t \in \mathrm{E} / \mathrm{R} \text { 且 } (\exists x) (x \in t \text { 且 } \mathrm{P} \{x \})
\]

被称为由P诱导的\{x\} 在过渡到商（关于x）关于R时。如果这个关系用P'表示 \{t\}，并且如果f是E到E/R的典范映射，那么关系

\[
y \in \mathrm{E} \text {   且   } \mathrm{P} ^ {\prime} \left\{f (y) \right\}
\]

（其中 \(y\) 中，而不出现在 \(P\{x\}\) )等价于( \(y \in E\) 并且 \(P\{y\}\) ），这一点可以立即验证。

